\documentclass[aspectratio=169]{beamer} 
\synctex=1
% 16:9 宽屏比例

\usepackage{../../Styles/mystyle-PPT}

\graphicspath{{./image/}}% 指定图片目录，后续可以直接使用图片文件名。

% beamer 主题、颜色与字体
\usetheme{metropolis}
\definecolor{mycolor}{RGB}{0, 176, 183}% 自定义青色
\definecolor{mycolor2}{RGB}{0, 191, 165}% 自定义亮绿色
% \usecolortheme[named=mycolor]{structure}% 让所有结构性元素（例如\item的序号等）都默认保持mycolor颜色

\usefonttheme{professionalfonts}% 让所有公式的字母、数字恢复为标准 LaTeX（Computer Modern）的衬线数学字体

\makeatother

% 演示文稿信息
\title{2016 年全国硕士研究生招生考试 \\ 数学一试题解析}
\subtitle{完整解答}
\author{邹文杰}
\institute{}
\date{\today}

\begin{document}

% 标题页
\begin{frame}
\titlepage
\end{frame}

% 目录
\begin{frame}{目录 CONTENTS}
\begin{columns}[T]
\column{0.48\textwidth}
\begin{block}{选择题}
1--8 小题
\end{block}
\begin{block}{填空题}
9--14 小题
\end{block}
\column{0.48\textwidth}
\begin{block}{解答题}
15--23 小题
\end{block}
\end{columns}
\end{frame}

% ============================================================
% 选择题部分
% ============================================================
\section{选择题}

\begin{frame}{选择题 (1--2)}
\begin{block}{第1题}
若反常积分 \(\displaystyle\int_{0}^{+\infty} \frac{1}{x^{a}(1+x)^{b}} \,\mathrm{d}x\) 收敛，则
$\mathrm{A.}\ a<1 \text{ 且 } b>1; \quad
\mathrm{B.}\ a>1 \text{ 且 } b>1; \quad
\mathrm{C.}\ a<1 \text{ 且 } a+b>1; \quad
\mathrm{D.}\ a>1 \text{ 且 } a+b>1$.
\\
\textcolor{red}{\textbf{答案：}$\mathbf{C}$} \\
解析：分段判敛，$x\to0^+$ 时要求 $a<1$，$x\to+\infty$ 时要求 $a+b>1$.
\end{block}
\end{frame}

\begin{frame}
\begin{block}{第2题}
已知函数 \(f(x)= \begin{cases}2(x-1), & x<1, \\ \ln x, & x \geq 1,\end{cases}\)
则 \(f(x)\) 的一个原函数是
\begin{enumerate}[A.]
\item \(F(x)=\begin{cases}(x-1)^{2}, & x<1, \\ x(\ln x-1), & x \geq 1;\end{cases}\)
\item \(F(x)=\begin{cases}(x-1)^{2}, & x<1, \\ x(\ln x+1)-1, & x \geq 1;\end{cases}\)
\item \(F(x)=\begin{cases}(x-1)^{2}, & x<1, \\ x(\ln x+1)+1, & x \geq 1;\end{cases}\)
\item \(F(x)=\begin{cases}(x-1)^{2}, & x<1, \\ x(\ln x-1)+1, & x \geq 1.\end{cases}\)
\end{enumerate}
\textcolor{red}{\textbf{答案：}$\mathbf{D}$} \\
解析：原函数在 $x=1$ 处连续，代入验证左右值相等即可.
\end{block}
\end{frame}

\begin{frame}{选择题 (3--4)}
\begin{block}{第3题}
若 \(y_1=(1+x^{2})^{2}-\sqrt{1+x^{2}}\)，\(y_2=(1+x^{2})^{2}+\sqrt{1+x^{2}}\) 是微分方程 \(y'+p(x) y=q(x)\) 的两个解，则 \(q(x)=\)
$\mathrm{A.}\ 3 x\left(1+x^{2}\right); \quad
\mathrm{B.}\ -3 x\left(1+x^{2}\right); \quad
\mathrm{C.}\ \dfrac{x}{1+x^{2}}; \quad
\mathrm{D.}\ -\dfrac{x}{1+x^{2}}$.
\\
\textcolor{red}{\textbf{答案：}$\mathbf{A}$} \\
解析：两解之差为齐次解，代入齐次方程求 $p(x)$，再代回原方程求 $q(x)$.
\end{block}
\end{frame}

\begin{frame}
\begin{block}{第4题}
已知函数
\[f(x)=\begin{cases} x, & x \leq 0, \\[4pt]
\dfrac{1}{n}, & \dfrac{1}{n+1}<x \leq \dfrac{1}{n},\ n=1,2,\cdots,\end{cases}\]
则
\begin{enumerate}[A.]
\item \(x=0\) 是 \(f(x)\) 的第一类间断点;
\item \(x=0\) 是 \(f(x)\) 的第二类间断点;
\item \(f(x)\) 在 \(x=0\) 处连续但不可导;
\item \(f(x)\) 在 \(x=0\) 处可导.
\end{enumerate}
\textcolor{red}{\textbf{答案：}$\mathbf{D}$} \\
解析：左右极限均为0，连续；左右导数均为1，故可导.
\end{block}
\end{frame}

\begin{frame}{选择题 (5--6)}
\begin{block}{第5题}
设矩阵 \(A\) 是可逆矩阵，且 \(A\) 与 \(B\) 相似，则下列结论错误的是
\begin{enumerate}[A.]
\item \(A^{T}\) 与 \(B^{T}\) 相似;
\item \(A^{-1}\) 与 \(B^{-1}\) 相似;
\item \(A+A^{T}\) 与 \(B+B^{T}\) 相似;
\item \(A+A^{-1}\) 与 \(B+B^{-1}\) 相似.
\end{enumerate}
\textcolor{red}{\textbf{答案：}$\mathbf{C}$} \\
解析：相似变换下转置对应 $P^T$ 而非 $P^{-1}$，故 $A+A^T$ 未必与 $B+B^T$ 相似.
\end{block}
\end{frame}

\begin{frame}
\begin{block}{第6题}
设二次型 \(f(x_{1}, x_{2}, x_{3})=x_{1}^{2}+x_{2}^{2}+x_{3}^{2}+4 x_{1} x_{2}+4 x_{1} x_{3}+4 x_{2} x_{3}\)，则 \(f(x_{1}, x_{2}, x_{3})=2\) 在空间直角坐标下表示的二次曲面为
\begin{enumerate}[A.]
\item 单叶双曲面;
\item 双叶双曲面;
\item 椭球面;
\item 柱面.
\end{enumerate}
\textcolor{red}{\textbf{答案：}$\mathbf{B}$} \\
解析：特征值为 $5,-1,-1$，标准型为 $5y_1^2-y_2^2-y_3^2=2$，对应双叶双曲面.
\end{block}
\end{frame}

\begin{frame}{选择题 (7--8)}
\begin{block}{第7题}
设随机变量 \(X \sim N(\mu, \sigma^{2})\ (\sigma>0)\)，记 \(p=P\{X \leq \mu+\sigma^{2}\}\)，则
\begin{enumerate}[A.]
\item \(p\) 随着 \(\mu\) 的增加而增加;
\item \(p\) 随着 \(\sigma\) 的增加而增加;
\item \(p\) 随着 \(\mu\) 的增加而减少;
\item \(p\) 随着 \(\sigma\) 的增加而减少.
\end{enumerate}
\textcolor{red}{\textbf{答案：}$\mathbf{B}$} \\
解析：标准化得 $p=\Phi(\sigma)$，$\Phi$ 单调递增，故 $p$ 随 $\sigma$ 增大而增大.
\end{block}
\end{frame}

\begin{frame}
\begin{block}{第8题}
随机试验 $E$ 有三种两两不相容的结果 \(A_{1},A_{2},A_{3}\)，且发生概率均为 $\dfrac13$. 将试验独立重复2次，$X$ 表示 $A_1$ 发生的次数，$Y$ 表示 $A_2$ 发生的次数，则 $\rho_{XY}=$
$\mathrm{A.}\ -\dfrac{1}{2} ;\quad
\mathrm{B.}\ -\dfrac{1}{3}; \quad
\mathrm{C.}\ \dfrac{1}{3}; \quad
\mathrm{D.}\ \dfrac{1}{2}$.
\\
\textcolor{red}{\textbf{答案：}$\mathbf{A}$} \\
解析：计算协方差与方差，得 $\operatorname{Cov}(X,Y)=-\dfrac29,\ D(X)=D(Y)=\dfrac49$，故相关系数为 $-\dfrac12$.
\end{block}
\end{frame}

% ============================================================
% 填空题部分
% ============================================================
\section{填空题}

\begin{frame}{填空题 (9--11)}
\begin{block}{第9题}
\(\displaystyle\lim _{x \to 0} \frac{\int_{0}^{x} t \ln (1+t \sin t) \,\mathrm{d} t}{1-\cos x^{2}}=\) \underline{$\quad\textcolor{red}{\dfrac{1}{2}}\quad$} 
\end{block}

\begin{block}{解析}
等价无穷小替换，$1-\cos x^2\sim \dfrac{x^4}{2}$，$\ln(1+t\sin t)\sim t^2$，分子等价于$\displaystyle\int_0^x t^3\mathrm{d}t=\frac{x^4}{4}$，故极限为$\dfrac{1}{2}$.
\end{block}
\end{frame}

\begin{frame}
\begin{block}{第10题}
向量场 \(\boldsymbol{A}(x, y, z)=(x+y+z)\boldsymbol{i}+xy\boldsymbol{j}+z\boldsymbol{k}\) 的旋度 \(\operatorname{rot}\boldsymbol{A}=\) \underline{$\quad\textcolor{red}{\boldsymbol{j}+(y-1)\boldsymbol{k}}\quad$} 
\end{block}

\begin{block}{解析}
由旋度公式 $\operatorname{rot}\boldsymbol{A}=\begin{vmatrix}\boldsymbol{i}&\boldsymbol{j}&\boldsymbol{k}\\[4pt]\dfrac{\partial}{\partial x}&\dfrac{\partial}{\partial y}&\dfrac{\partial}{\partial z}\\[4pt]P&Q&R\end{vmatrix}$，代入计算得分量分别为 $0,\,1,\,y-1$.
\end{block}
\end{frame}

\begin{frame}
\begin{block}{第11题}
设函数 \(z=z(x, y)\) 由方程 \((x+1) z-y^{2}=x^{2} f(x-z, y)\) 确定，则 \(\left.\mathrm{d} z\right|_{(0,1)}=\) \underline{$\quad \textcolor{red}{-\mathrm{d} x+2\mathrm{d} y}\quad$} 
\end{block}

\begin{block}{解析}
代入 $(0,1)$ 得 $z=1$；两边求微分并代入该点坐标，解得 $\dfrac{\partial z}{\partial x}=-1$，$\dfrac{\partial z}{\partial y}=2$，故 $\mathrm{d}z=-\mathrm{d}x+2\mathrm{d}y$.
\end{block}
\end{frame}

\begin{frame}{填空题 (12--14)}
\begin{block}{第12题}
设函数 \(f(x)=\arctan x-\dfrac{x}{1+a x^{2}}\)，若 \(f^{\prime \prime \prime}(0)=1\)，则 \(a=\) \underline{$\quad \textcolor{red}{\dfrac{1}{2}}\quad$} 
\end{block}

\begin{block}{解析}
利用麦克劳林展开，$\arctan x=x-\dfrac{x^3}{3}+o(x^3)$，$\dfrac{x}{1+ax^2}=x-ax^3+o(x^3)$，故 $f(x)=\left(a-\dfrac{1}{3}\right)x^3+o(x^3)$，由 $f'''(0)=6\left(a-\dfrac{1}{3}\right)=1$ 得 $a=\dfrac{1}{2}$.
\end{block}
\end{frame}

\begin{frame}
\begin{block}{第13题}
行列式 \(\begin{vmatrix}
\lambda & -1 & 0 & 0 \\
0 & \lambda & -1 & 0 \\
0 & 0 & \lambda & -1 \\
4 & 3 & 2 & \lambda+1
\end{vmatrix} =\) \underline{$\quad \textcolor{red}{\lambda^{4}+\lambda^{3}+2\lambda^{2}+3\lambda+4} \quad$} 
\end{block}

\begin{block}{解析}
按第一列展开递推，逐级计算低阶行列式，利用三对角行列式的展开规律化简即得结果.
\end{block}
\end{frame}

\begin{frame}
\begin{block}{第14题}
设 \(X_{1}, X_{2}, \cdots, X_{n}\) 为来自总体 \(N(\mu, \sigma^{2})\) 的简单随机样本，样本均值 \(\bar{X}=9.5\)，$\mu$ 的置信度0.95双侧置信上限为10.8，则置信区间为 \underline{$\quad \textcolor{red}{(8.2,\ 10.8)}\quad$} 
\end{block}

\begin{block}{解析}
正态总体双侧置信区间关于样本均值对称，下限为 $9.5-(10.8-9.5)=8.2$，故置信区间为 $(8.2,\,10.8)$.
\end{block}
\end{frame}

% ============================================================
% 解答题部分
% ============================================================
\section{解答题}

\begin{frame}{第15题}
\begin{block}{题目}
已知平面区域 \(D=\{(r, \theta) \mid 2 \leqslant r \leqslant 2(1+\cos \theta),\ -\dfrac{\pi}{2} \leqslant \theta \leqslant \dfrac{\pi}{2}\}\)，计算二重积分 \(\displaystyle\iint_{D} x\,\mathrm{d}\sigma\).
\end{block}
\begin{block}{解答}
极坐标下 \(x=r\cos\theta,\ \mathrm{d}\sigma=r\mathrm{d}r\mathrm{d}\theta\)，利用对称性
\begin{align*}
\iint_D x\,\mathrm{d}\sigma &= \int_{-\frac{\pi}{2}}^{\frac{\pi}{2}} \cos\theta \,\mathrm{d}\theta \int_{2}^{2(1+\cos\theta)} r^2 \,\mathrm{d}r \\
&= \frac{16}{3}\int_{0}^{\frac{\pi}{2}} \cos\theta\big[(1+\cos\theta)^3-1\big]\mathrm{d}\theta \\
&= \frac{16}{3}\int_{0}^{\frac{\pi}{2}} \big(3\cos^2\theta+3\cos^3\theta+\cos^4\theta\big)\mathrm{d}\theta \\
&= 5\pi+\frac{32}{3}.
\end{align*}
\end{block}
\end{frame}

\begin{frame}{第16题}
\begin{block}{题目}
设函数 \(y(x)\) 满足方程 \(y''+2 y'+k y=0\)，其中 \(0<k<1\).
\begin{enumerate}[(1)]
\item 证明：反常积分 \(\displaystyle\int_{0}^{+\infty} y(x) \,\mathrm{d}x\) 收敛；

\item 若 \(y(0)=1,\ y'(0)=1\)，求 \(\displaystyle\int_{0}^{+\infty} y(x) \,\mathrm{d}x\) 的值.
\end{enumerate} 
\end{block}
\end{frame}

\begin{frame}
\begin{block}{解答}
\begin{enumerate}[(1)]
\item 特征方程 \(r^2+2r+k=0\)，根为 \(r=-1\pm\sqrt{1-k}\)，均为负实根，故 \(y(x)\to0,\ y'(x)\to0\ (x\to+\infty)\)，积分收敛.

\item 方程两边在 \([0,+\infty)\) 积分
\begin{align*}
\int_0^{+\infty} y''\mathrm{d}x + 2\int_0^{+\infty} y'\mathrm{d}x + k\int_0^{+\infty} y\mathrm{d}x = 0,
\end{align*}
代入边界值得 \(-1-2 + kI=0\)，故 \(I=\dfrac{3}{k}\).
\end{enumerate}    
\end{block}
\end{frame}

\begin{frame}{第17题}
\begin{block}{题目}
设函数 \(f(x, y)\) 满足 \(\dfrac{\partial f(x, y)}{\partial x}=(2 x+1) e^{2 x-y}\)，且 \(f(0, y)=y+1\)，\(L_t\) 是从 \((0,0)\) 到 \((t,0)\) 再到 \((t,1)\) 的折线段. 计算曲线积分 \(I(t)=\displaystyle\int_{L_{t}} \frac{\partial f}{\partial x} \mathrm{d} x+\frac{\partial f}{\partial y} \mathrm{d} y\)，并求 \(I(t)\) 的最小值.
\end{block}
\begin{block}{解答}
对 $x$ 积分得 \(f(x,y)=x e^{2x-y}+y+1\)，积分与路径无关，故
\begin{align*}
I(t)=f(t,1)-f(0,0)=t e^{2t-1}+1.
\end{align*}
求导得 \(I'(t)=(2t+1)e^{2t-1}\)，令 \(I'(t)=0\) 得 \(t=-\dfrac12\)，
最小值为 \(I\!\left(-\dfrac12\right)=1-\dfrac{1}{2e^2}\).
\end{block}
\end{frame}

\begin{frame}{第18题}
\begin{block}{题目}
设有界区域 \(\Omega\) 由平面 \(2 x+y+2 z=2\) 与三个坐标平面围成，\(\Sigma\) 为 \(\Omega\) 整个表面的外侧，计算曲面积分
\[I=\iint_{\Sigma}\left(x^{2}+1\right) \mathrm{d} y \mathrm{d} z-2 y \,\mathrm{d} z \mathrm{d} x+3 z \,\mathrm{d} x \mathrm{d} y.\]
\end{block}
\begin{block}{解答}
由高斯公式，\(\operatorname{div}\boldsymbol{F}=2x+1\)，故
\begin{align*}
I=\iiint_\Omega (2x+1)\,\mathrm{d}V.
\end{align*}
计算三重积分：体积 \(\displaystyle V=\frac13\)，\(\displaystyle\iiint_\Omega 2x\,\mathrm{d}V=\frac16\)，
因此 \(I=\dfrac16+\dfrac13=\dfrac12\).
\end{block}
\end{frame}

\begin{frame}{第19题}
\begin{block}{题目}
已知函数 \(f(x)\) 可导，且 \(f(0)=1,\ 0<f'(x)<\dfrac{1}{2}\). 设数列 \(\{x_{n}\}\) 满足 \(x_{n+1}=f(x_{n})\ (n=1,2,\cdots)\). 证明：
(1) 级数 \(\displaystyle\sum_{n=1}^{\infty}(x_{n+1}-x_{n})\) 绝对收敛；
(2) \(\displaystyle\lim_{n \to \infty} x_{n}\) 存在，且 \(0<\lim_{n \to \infty} x_{n}<2\).
\end{block}
\begin{block}{证明}
(1) 由拉格朗日中值定理
\[|x_{n+1}-x_n|=|f(x_n)-f(x_{n-1})|=|f'(\xi)|\cdot|x_n-x_{n-1}|<\frac12|x_n-x_{n-1}|,\]
由比值判别法知级数绝对收敛.

(2) 级数绝对收敛故 \(\{x_n\}\) 收敛，设极限为 $a$，则 $a=f(a)$.
令 $g(x)=f(x)-x$，$g(0)=1>0,\ g(2)<0$，由零点定理与单调性得 $0<a<2$.
\end{block}
\end{frame}

\begin{frame}{第20题}
\begin{block}{题目}
设矩阵 \(A=\begin{pmatrix}1&-1&-1\\2&a&1\\-1&1&a\end{pmatrix}\)，\(B=\begin{pmatrix}2&2\\1&a\\-a-1&-2\end{pmatrix}\). 当 $a$ 为何值时，方程 \(AX=B\) 无解、有唯一解、有无穷多解？在有解时求解.
\end{block}
\begin{block}{解答}
对增广矩阵做初等行变换：
\begin{enumerate}[(1)]
\item 当\(a=-2\) 时，秩不等，无解；

\item 当\(a\neq-2\) 且 \(a\neq1\) 时，秩为3，有唯一解
\[X=\begin{pmatrix}1 & \dfrac{3a}{a+2}\\[4pt] 0 & \dfrac{a-4}{a+2}\\[4pt] -1 & 0\end{pmatrix};\]
\end{enumerate}
\end{block}
\end{frame}
\begin{frame}
\begin{enumerate}[(3)]
\item 当\(a=1\) 时，有无穷多解
\[X=\begin{pmatrix}1 & 1\\ -1 & -1\\ 0 & 0\end{pmatrix}+\begin{pmatrix}0 & 0\\ -c_1 & -c_2\\ c_1 & c_2\end{pmatrix},\quad c_1,c_2\in\mathbb{R}.\]
\end{enumerate}
\end{frame}

\begin{frame}{第21题}
\begin{block}{题目}
已知矩阵 \(A=\begin{pmatrix}0 & -1 & 1 \\ 2 & -3 & 0 \\ 0 & 0 & 0\end{pmatrix}\).
\begin{enumerate}[(1)]
\item 求 \(A^{99}\)；
\item 设3阶矩阵 \(B=(\alpha_1,\alpha_2,\alpha_3)\) 满足 \(B^2=BA\)，记 \(B^{100}=(\beta_1,\beta_2,\beta_3)\)，将 \(\beta_1,\beta_2,\beta_3\) 表示为 \(\alpha_1,\alpha_2,\alpha_3\) 的线性组合.
\end{enumerate}
\end{block}
\end{frame}

\begin{frame}
\begin{block}{解答}
\begin{enumerate}[(1)]
\item $A$ 的特征值为 $0,-1,-2$，对角化后计算得
\[A^{99}=\begin{pmatrix}
2^{99}-2 & 1-2^{99} & 2-2^{98} \\
2^{100}-2 & 1-2^{100} & 2-2^{99} \\
0 & 0 & 0
\end{pmatrix}.\]

\item 由 \(B^{100}=B\cdot A^{99}\)，得
\begin{align*}
\beta_1&=(2^{99}-2)\alpha_1+(2^{100}-2)\alpha_2, \\
\beta_2&=(1-2^{99})\alpha_1+(1-2^{100})\alpha_2, \\
\beta_3&=(2-2^{98})\alpha_1+(2-2^{99})\alpha_2.
\end{align*}
\end{enumerate}
\end{block}
\end{frame}

\begin{frame}{第22题}
\begin{block}{题目}
设二维随机变量 \((X, Y)\) 在区域 \(D=\{(x, y) \mid 0<x<1,\ x^{2}<y<\sqrt{x}\}\) 上服从均匀分布，令
\[U=\begin{cases} 1, & X \leq Y, \\ 0, & X>Y. \end{cases}\]
\begin{enumerate}[(1)]
\item 写出 \((X, Y)\) 的概率密度；

\item 问 \(U\) 与 \(X\) 是否相互独立，并说明理由；

\item 求 \(Z=U+X\) 的分布函数 \(F(z)\).
\end{enumerate}
\end{block}
\end{frame}

\begin{frame}
\begin{block}{解答}
\begin{enumerate}[(1)]
\item 区域面积为 \(\dfrac13\)，故 \(f(x,y)=\begin{cases}3, & (x,y)\in D,\\ 0, & \text{其他}.\end{cases}\)

\item 不独立，可验证 \(P(U=0,X\leqslant\tfrac12)\neq P(U=0)P(X\leqslant\tfrac12)\).

\item 分段计算得
\[F(z)= \begin{cases}
0, & z \leq 0, \\[4pt]
\dfrac{3}{2} z^{2}-z^{3}, & 0<z \leq 1, \\[4pt]
2(z-1)^{\frac{3}{2}}-\dfrac{3}{2} z^{2}+3 z-1, & 1<z \leq 2, \\[4pt]
1, & z>2.
\end{cases}\]
\end{enumerate}
\end{block}
\end{frame}

\begin{frame}{第23题}
\begin{block}{题目}
设总体 \(X\) 的概率密度为
\[f(x ; \theta)=\begin{cases}
\dfrac{3 x^{2}}{\theta^{3}}, & 0<x<\theta, \\[4pt]
0, & \text{其他},
\end{cases}\]
其中 \(\theta \in(0,+\infty)\) 为未知参数. \(X_{1},X_{2},X_{3}\) 为来自总体 \(X\) 的简单随机样本，令 \(T=\max \left\{X_{1}, X_{2}, X_{3}\right\}\).
\begin{enumerate}[(1)]
\item 求 \(T\) 的概率密度；
\item 确定 \(a\)，使得 \(aT\) 为 \(\theta\) 的无偏估计.
\end{enumerate}
\end{block}
\end{frame}

\begin{frame}
\begin{block}{解答}
\begin{enumerate}[(1)]
\item 先求分布函数 \(F_T(t)=[F_X(t)]^3\)，求导得
\[f_T(t)=\begin{cases}
\dfrac{9 t^{8}}{\theta^{9}}, & 0<t<\theta, \\[4pt]
0, & \text{其他}.
\end{cases}\]

\item 计算期望 \(E(T)=\dfrac{9}{10}\theta\)，令 \(E(aT)=\theta\)，得 \(a=\dfrac{10}{9}\).
\end{enumerate}
\end{block}
\end{frame}

\begin{frame}
\centering
\vfill
\Huge \textcolor{mycolor}{感谢观看} \\
\vfill
\end{frame}

\end{document}